% Part: turing-machines
% Chapter: machines-computations
% Section: representing-tms

\documentclass[../../../include/open-logic-section]{subfiles}

\begin{document}

\olfileid{tur}{mac}{rep}
\olsection{ટ્યુરિંગ મશીનોનું નિરૂપણ}

\begin{explain}
ટ્યુરિંગ મશીનોને \emph{અવસ્થા-આલેખો} વડે ચિત્રરૂપે દર્શાવી
શકાય છે. આવા આલેખોમાં તીરોથી જોડાયેલાં અવસ્થાનાં ખાનાં હોય
છે. દરેક અવસ્થાનું ખાનું મશીનની એક અવસ્થા દર્શાવે છે. દરેક
તીર તે અવસ્થામાંથી અમલમાં મૂકી શકાય એવી એક સૂચના દર્શાવે છે;
તે સૂચનાની વિગતો સંબંધિત તીરની ઉપર અથવા નીચે લખેલી હોય
છે. નીચેના મશીનમાં માત્ર બે આંતરિક અવસ્થાઓ, $q_0$ અને
$q_1$, તથા એક સૂચના છે:
\[
\begin{tikzpicture}[->,>=stealth',shorten >=1pt,auto,node distance=2.8cm,
                    semithick]
  \tikzstyle{every state}=[fill=none,draw=black,text=black]

  \node[initial,state]   (A)              {$q_0$};
  \node[state]   (B) [right of=A] {$q_1$};

  \path (A) edge  node {\TMtrans{\TMblank}{\TMstroke}{\TMright}} (B);
\end{tikzpicture}
\]
યાદ કરો કે ટ્યુરિંગ મશીન પાસે વાંચન-લેખન હેડ અને તેના પર
ઇનપુટ લખેલી ટેપ હોય છે. આ સૂચનાનો અર્થ છે: \emph{અવસ્થા
$q_0$માં $\TMblank$ વાંચતા હો, તો $\TMstroke$ લખો, જમણે ખસો
અને અવસ્થા $q_1$માં જાઓ}. આ વાત સંક્રમણ વિધેયમાં
$\tuple{q_0, \TMblank}$ને $\tuple{q_1, \TMstroke,
\TMright}$ સાથે જોડવા સમકક્ષ છે.
\end{explain}

\begin{ex}
\emph{સમ મશીન}: ટેપ પર $\TMstroke$ની સંખ્યા સમ હોય ત્યારે
અને ત્યારે જ નીચેનું ટ્યુરિંગ મશીન થંભે છે (એવી ધારણા હેઠળ કે
ટેપ પરના બધા $\TMstroke$ પ્રથમ $\TMblank$ પહેલાં આવે છે).
\[
\begin{tikzpicture}[->,>=stealth',shorten >=1pt,auto,node distance=2.8cm,
                    semithick]
  \tikzstyle{every state}=[fill=none,draw=black,text=black]

  \node[initial,state]         (A)              {$q_0$};
  \node[state]         (B) [right of=A] {$q_1$};

  \path (A) edge [bend left] node {\TMtrans{\TMstroke}{\TMstroke}{\TMright}} (B)
        (B) edge [loop above] node {\TMtrans{\TMblank}{\TMblank}{\TMright}} (B)
            edge [bend left] node {\TMtrans{\TMstroke}{\TMstroke}{\TMright}} (A);
\end{tikzpicture}
\]

આ અવસ્થા-આલેખને નીચેનું સંક્રમણ વિધેય અનુરૂપ છે:
\begin{align*}
\delta(q_0, \TMstroke) & = \tuple{q_1, \TMstroke, \TMright},\\
\delta(q_1, \TMstroke) & = \tuple{q_0, \TMstroke, \TMright},\\
\delta(q_1, \TMblank)  & = \tuple{q_1, \TMblank, \TMright}
\end{align*}
\end{ex}

\begin{explain}
ઉપરનું મશીન ઇનપુટમાં સ્ટ્રોકની સંખ્યા સમ હોય ત્યારે
જ થંભે છે. નહીંતર મશીન (સૈદ્ધાંતિક રીતે) અનંત સમય સુધી
ચાલુ રહે છે. કોઈ પણ મશીન અને ઇનપુટ માટે આઉટપુટ નક્કી
કરવા મશીનની \emph{સંરૂપણો}નો ક્રમ અનુસરવો શક્ય છે.
સંરૂપણની ઔપચારિક વ્યાખ્યા આપણે આગળ આપીશું. હાલ પૂરતું,
મશીન ચાલે ત્યારે જુદા જુદા સમયે તેની સ્થિતિ દર્શાવતા આલેખોની
શ્રેણી તરીકે સંરૂપણોને સહજ રીતે સમજી શકાય. સંરૂપણમાં ટેપની
સામગ્રી, મશીનની અવસ્થા અને વાંચન-લેખન હેડનું સ્થાન દેખાય છે.

ચાર $\TMstroke$નું ઇનપુટ આપીએ ત્યારે સમ મશીનનાં સંરૂપણો
ક્રમશઃ જોઈએ. આ કિસ્સામાં મશીન થંભશે એવી આપણી અપેક્ષા
છે. પછી તેને ત્રણ $\TMstroke$નું ઇનપુટ આપીશું; ત્યારે તે
હંમેશાં ચાલતું રહેશે.

મશીન અવસ્થા~$q_0$માં શરૂ થાય છે અને સૌથી ડાબો
$\TMstroke$ વાંચે છે. મશીનની આ પ્રારંભિક સ્થિતિ આ રીતે
દર્શાવી શકાય:
\[
\TMendtape \TMstroke_0 \TMstroke \TMstroke \TMstroke \TMblank \ldots
\]
ઉપરનું સંરૂપણ સમજવું સહેલું છે. મશીન પ્રથમ અવસ્થામાં
સૌથી ડાબો $\TMstroke$ વાંચતાં શરૂ થાય છે. પ્રથમ
$\TMstroke$ને અવસ્થાના નામનો નિમ્નસૂચક લગાવીને આ દર્શાવ્યું
છે. આ સમયે લાગુ પડતી સૂચના
$\delta(q_0, \TMstroke) = \tuple{q_1, \TMstroke, \TMright}$
છે. તેથી મશીન ટેપ પર જમણે ખસે છે અને અવસ્થા~$q_1$માં
જાય છે.
\[
\TMendtape \TMstroke \TMstroke_1 \TMstroke \TMstroke \TMblank \ldots
\]
હવે મશીન અવસ્થા~$q_1$માં $\TMstroke$ વાંચતું હોવાથી,
આપણે સૂચના $\delta(q_1, \TMstroke) = \tuple{q_0,
  \TMstroke, \TMright}$ને ``અનુસરવી'' પડશે. પરિણામે આ
સંરૂપણ મળે છે:
\[
\TMendtape \TMstroke \TMstroke \TMstroke_0 \TMstroke \TMblank \ldots
\]
મશીન આગળ ચાલે છે ત્યારે આ જ ક્રમમાં નિયમો ફરી લાગુ પડે
છે અને નીચેના બે સંરૂપણો મળે છે:
\[
\TMendtape \TMstroke \TMstroke \TMstroke \TMstroke_1 \TMblank \ldots
\]
\[
\TMendtape \TMstroke \TMstroke \TMstroke \TMstroke \TMblank_0 \ldots
\]
હવે મશીન અવસ્થા~$q_0$માં $\TMblank$ વાંચે છે. સંક્રમણ
આલેખ પરથી સહેલાઈથી દેખાય છે કે હવે અમલમાં મૂકવા માટે કોઈ
સૂચના નથી, એટલે મશીન થંભી ગયું છે. તેનો અર્થ છે કે
ઇનપુટ સ્વીકારાયું છે.

હવે ધારો કે મશીનને ત્રણ $\TMstroke$નું ઇનપુટ આપીએ.
ટેપ પરના ઇનપુટમાં નાનો તફાવત હોવા છતાં એ જ સૂચનાઓ લાગુ
પડતી હોવાથી શરૂઆતનાં કેટલાંક સંરૂપણો સરખાં હોય છે:
\[
\TMendtape \TMstroke_0 \TMstroke \TMstroke \TMblank \ldots
\]
\[
\TMendtape \TMstroke \TMstroke_1 \TMstroke \TMblank \ldots
\]
\[
\TMendtape \TMstroke \TMstroke \TMstroke_0 \TMblank \ldots
\]
\[
\TMendtape \TMstroke \TMstroke \TMstroke \TMblank_1 \ldots
\]
મશીન હવે બધા $\TMstroke$ વટાવી ચૂક્યું છે અને
અવસ્થા~$q_1$માં $\TMblank$ વાંચે છે. આલેખમાં બતાવ્યા
પ્રમાણે, અહીં $\delta(q_1, \TMblank) =\tuple{q_1,
\TMblank, \TMright}$ સ્વરૂપની સૂચના છે. ટેપ જમણી તરફ
અનંત સુધી $\TMblank$થી ભરેલી હોવાથી, મશીન
અવસ્થા~$q_1$માં રહીને વધુ ને વધુ જમણે ખસતું આ સૂચના
\emph{હંમેશાં} ચલાવશે. મશીન ક્યારેય થંભશે નહીં અને ઇનપુટ
સ્વીકારશે નહીં.
\end{explain}

\begin{explain}
દરેક મશીન થંભતું નથી, એ નોંધવું અગત્યનું છે. થંભવાનો અર્થ
અમલમાં મૂકવા માટેની સૂચનાઓ ખૂટી જવી એવો હોય, તો દરેક
અવસ્થામાં દરેક પ્રતીક માટે બહાર જતું તીર છે તેની ખાતરી
કરીને ક્યારેય ન થંભે એવું મશીન બનાવી શકાય. અવસ્થા~$q_0$માં
$\TMblank$ વાંચવા માટેની સૂચના ઉમેરીને સમ મશીનને અનંત સમય
સુધી ચાલે એવું બનાવી શકાય.
\end{explain}

\begin{ex}
\[
\begin{tikzpicture}[->,>=stealth',shorten >=1pt,auto,node distance=2.8cm,
                    semithick]
  \tikzstyle{every state}=[fill=none,draw=black,text=black]

  \node[initial,state] (A)              {$q_0$};
  \node[state]         (B) [right of=A] {$q_1$};

  \path
  (A) edge [bend left] node {\TMtrans{\TMstroke}{\TMstroke}{\TMright}} (B)
      edge [loop above] node {\TMtrans{\TMblank}{\TMblank}{\TMright}} (A)
  (B) edge [loop above] node {\TMtrans{\TMblank}{\TMblank}{\TMright}} (B)
      edge [bend left] node {\TMtrans{\TMstroke}{\TMstroke}{\TMright}} (A);
\end{tikzpicture}
\]
\end{ex}

\begin{explain}
ટ્યુરિંગ મશીનો દર્શાવવાની બીજી રીત \emph{મશીન-કોષ્ટકો}
છે. કોષ્ટકની $x$-અક્ષ પર ટેપની વર્ણમાળા અને $y$-અક્ષ પર
મશીનની અવસ્થાઓનો ગણ દર્શાવાય છે. દરેક અવસ્થા અને પ્રતીકના
છેદે સૂચનાનો બાકીનો ભાગ---નવી અવસ્થા, નવું પ્રતીક અને ખસવાની
દિશા---લખાય છે. કઈ અવસ્થામાં અને કયું પ્રતીક વાંચતાં મશીન
થંભે છે તે કોષ્ટકથી સહેલાઈથી જાણી શકાય છે. કોષ્ટકમાં કોઈ
ખાનું ખાલી હોય, તો ત્યાં મશીન થંભી શકે છે. અવસ્થા-આલેખો
અને સૂચનાના ગણમાં મશીન ક્યાં થંભે છે તે હંમેશાં તરત દેખાતું
નથી; મશીન-કોષ્ટકમાં ખાલી ખાનાં જોઈને આ બિંદુઓ તરત ઓળખી
શકાય છે.
\end{explain}

\begin{ex}
સમ મશીનનું મશીન-કોષ્ટક આ છે:
\[
\centering
\begin{tabular}{lllll}
\cline{2-4}
\multicolumn{1}{l|}{}      & \multicolumn{1}{c|}{$\TMblank$}
& \multicolumn{1}{c|}{$\TMstroke$}     & \multicolumn{1}{c|}{$\TMendtape$}   &  \\ \cline{1-4}
\multicolumn{1}{|l|}{$q_0$} & \multicolumn{1}{l|}{}
& \multicolumn{1}{l|}{\TMtrans{\TMstroke}{q_1}{\TMright}} &
\multicolumn{1}{l|}{\phantom{\TMtrans{\TMstroke}{q_1}{\TMright}}} &  \\ \cline{1-4}
\multicolumn{1}{|l|}{$q_1$} & \multicolumn{1}{l|}{\TMtrans{\TMblank}{q_1}{\TMright}}
& \multicolumn{1}{l|}{\TMtrans{\TMstroke}{q_0}{\TMright}} &
\multicolumn{1}{l|}{\phantom{\TMtrans{\TMstroke}{q_1}{\TMright}}} &  \\ \cline{1-4}
\end{tabular}
\]
જેમ જોઈ શકાય છે, મશીન અવસ્થા~$q_0$માં $\TMblank$
વાંચે ત્યારે થંભે છે.
\end{ex}

\begin{explain}
અત્યાર સુધી આપણે ઇનપુટ વાંચીને સ્વીકારતાં મશીનો જ જોયાં
છે. પરંતુ ટ્યુરિંગ મશીન વાંચી પણ શકે છે અને લખી પણ શકે
છે. આવા મશીનનું એક ઉદાહરણ (જોકે આવાં ઘણાં ઉદાહરણો છે)
\emph{બમણું કરનાર મશીન} છે. ટેપ પર $n$ $\TMstroke$નો
એક સળંગ સમૂહ આપીને શરૂ કરીએ, તો તે $2n$ $\TMstroke$નો
સમૂહ આઉટપુટ તરીકે આપે છે.
\end{explain}

\begin{ex}
  \ollabel{ex:doubler} બમણું કરનાર મશીન બનાવતાં પહેલાં સમસ્યા
  ઉકેલવાની \emph{રીત} નક્કી કરવી જરૂરી છે. મશીન (આપણે તેને
  જેમ ઘડ્યું છે તે મુજબ) કેટલા $\TMstroke$ વાંચ્યા છે તે યાદ
  રાખી શકતું નથી. તેથી ટેપ પરના બધા $\TMstroke$નો હિસાબ
  રાખવાની રીત શોધવી પડશે. એક રીત ઇનપુટને આઉટપુટથી
  એક $\TMblank$ વડે અલગ કરવાની છે. પછી મશીન ઇનપુટમાંથી
  પ્રથમ $\TMstroke$ ભૂંસી શકે, બાકીના ઇનપુટને વટાવી શકે,
  એક $\TMblank$ છોડી શકે અને બે નવા $\TMstroke$ લખી
  શકે. ત્યાર પછી મશીન પાછું જઈ ઇનપુટનો બીજો $\TMstroke$
  શોધશે અને તેને પણ બમણો કરશે. ઇનપુટના દરેક એક
  $\TMstroke$ માટે તે આઉટપુટમાં બે $\TMstroke$ લખશે.
  મશીન આગળ વધે તેમ ઇનપુટ ભૂંસવાથી કોઈ $\TMstroke$ છૂટી
  ન જાય કે બે વાર બમણો ન થાય તેની ખાતરી કરી શકાય. આખું
  ઇનપુટ ભૂંસાઈ જાય ત્યારે ટેપ પર $2n$ $\TMstroke$ બાકી
  હશે. મળેલા ટ્યુરિંગ મશીનનો અવસ્થા-આલેખ
  \olref{fig:doubler}માં દર્શાવ્યો છે.
  \begin{figure}
\[
\begin{tikzpicture}[->,>=stealth',shorten >=1pt,auto,node distance=2.8cm,
                    semithick]
  \tikzstyle{every state}=[fill=none,draw=black,text=black]

  \node[initial,state] (1)              {$q_0$};
  \node[state]         (2) [right of=1] {$q_1$};
  \node[state]         (3) [right of=2] {$q_2$};
  \node[state]         (4) [below of=3] {$q_3$};
  \node[state]         (5) [left of=4]  {$q_4$};
  \node[state]         (6) [left of=5]  {$q_5$};

  \path (1) edge node {\TMtrans{\TMstroke}{\TMblank}{\TMright}} (2)
    (2) edge [loop above] node {\TMtrans{\TMstroke}{\TMstroke}{\TMright}} (2)
      edge node {\TMtrans{\TMblank}{\TMblank}{\TMright}} (3)
    (3) edge [loop above] node {\TMtrans{\TMstroke}{\TMstroke}{\TMright}} (3)
        edge node {\TMtrans{\TMblank}{\TMstroke}{\TMright}} (4)
    (4) edge [loop below] node {\TMtrans{\TMblank}{\TMstroke}{\TMleft}} (4)
        edge node {\TMtrans{\TMstroke}{\TMstroke}{\TMleft}} (5)
    (5) edge [loop below]  node {\TMtrans{\TMstroke}{\TMstroke}{\TMleft}} (5)
        edge              node {\TMtrans{\TMblank}{\TMblank}{\TMleft}} (6)
    (6) edge [loop below] node {\TMtrans{\TMstroke}{\TMstroke}{\TMleft}} (6)
        edge              node {\TMtrans{\TMblank}{\TMblank}{\TMright}} (1);
\end{tikzpicture}
\]
\caption{બમણું કરનાર મશીન}
\ollabel{fig:doubler}
\end{figure}
\end{ex}

\begin{prob}
કોઈ પણ એક ઇનપુટ પસંદ કરો અને \olref[tur][mac][rep]{ex:doubler}માંના
બમણું કરનાર મશીનનાં સંરૂપણોનો ક્રમ અનુસરો.
\end{prob}

\begin{prob}
$\{\TMendtape,\TMblank, A, B\}$ વર્ણમાળા ધરાવતું એવું
ટ્યુરિંગ મશીન રચો કે જેમાં $A$ અને $B$ની સંખ્યા સરખી
હોય \emph{અને} બધા $A$, બધા $B$ પહેલાં આવતા હોય,
એવી $A$ અને $B$ની કોઈ પણ શબ્દમાળા તે સ્વીકારે,
એટલે તેના પર થંભે; અને $A$ની સંખ્યા $B$ની સંખ્યા
જેટલી ન હોય અથવા બધા $A$ બધા $B$ પહેલાં ન આવતા
હોય એવી શબ્દમાળા તે નકારે, એટલે તેના પર ન થંભે.
(દાખલા તરીકે, મશીને $AABB$ અને $AAABBB$ સ્વીકારવાં,
પણ $AAB$ અને $AABBAABB$ બંને નકારવાં જોઈએ.)
\end{prob}

\begin{prob}
$\{\TMendtape,\TMblank, A, B\}$ વર્ણમાળા ધરાવતું એવું
ટ્યુરિંગ મશીન રચો કે જે $A$ અને $B$ની કોઈ પણ
શબ્દમાળા~$\alpha$ ઇનપુટ તરીકે લે અને તેની નકલ કરીને
$\alpha\alpha$ સ્વરૂપનું આઉટપુટ આપે. (દાખલા તરીકે,
ઇનપુટ $ABBA$નું આઉટપુટ $ABBAABBA$ થવું જોઈએ.)
\end{prob}

\begin{prob}
\emph{વર્ણક્રમમાં?:} $\{\TMendtape,\TMblank, A, B\}$ વર્ણમાળા
ધરાવતું એવું ટ્યુરિંગ મશીન રચો કે જેને $A$ અને $B$નો
સાન્ત અનુક્રમ ઇનપુટ તરીકે આપતાં બધા $A$, બધા $B$ની
ડાબે આવે છે કે નહીં તે તપાસે. મશીને ઇનપુટ શબ્દમાળા ટેપ
પર જ રાખવી જોઈએ અને શબ્દમાળા ``વર્ણક્રમમાં'' હોય તો
થંભવું, નહીંતર હંમેશાં ચક્રમાં ચાલવું જોઈએ.
\end{prob}

\begin{prob}
\emph{વર્ણક્રમમાં ગોઠવનાર:} $\{\TMendtape,\TMblank, A, B\}$
વર્ણમાળા ધરાવતું એવું ટ્યુરિંગ મશીન રચો કે જે $A$ અને
$B$ના સાન્ત અનુક્રમને ઇનપુટ તરીકે લઈ તેમને ફરી ગોઠવે,
જેથી બધા $A$, બધા~$B$ની ડાબે હોય. (દાખલા તરીકે,
$BABAA$ અનુક્રમ $AAABB$ બનવો જોઈએ અને $ABBABB$
અનુક્રમ $AABBBB$ બનવો જોઈએ.)
\end{prob}

\end{document}
